% ============================================================================
% CHAPTER 3: METHODOLOGY
% ============================================================================

\chapter{METHODOLOGY}

\section{SYSTEM OVERVIEW}

PassiveGuard follows a modular three-tier architecture comprising a frontend SDK for data collection, a backend API for processing, and a machine learning model for classification. The system is designed for minimal latency, maximum privacy, and easy integration.

\subsection{Design Principles}

\begin{itemize}
    \item \textbf{Privacy by Design:} Fingerprints are hashed before transmission; no PII is collected
    \item \textbf{Passive Collection:} All data gathering occurs transparently without user interaction
    \item \textbf{Defense in Depth:} Multiple signal categories provide layered detection
    \item \textbf{Minimal Footprint:} Lightweight SDK with minimal performance impact
    \item \textbf{Pluggable Architecture:} Easy integration with any web application
\end{itemize}

\subsection{Workflow Overview}

The complete bot detection workflow:
\begin{enumerate}
    \item \textbf{Initialization:} SDK loads and begins passive data collection
    \item \textbf{Environmental Collection:} Browser, device, and fingerprint data captured
    \item \textbf{Behavioral Monitoring:} Mouse, keyboard, scroll, and timing patterns tracked
    \item \textbf{Verification Request:} Collected data sent to backend API
    \item \textbf{Feature Extraction:} Raw data transformed into 56 ML features
    \item \textbf{Model Inference:} XGBoost classifier predicts human vs bot
    \item \textbf{Result Generation:} Verification result with confidence score returned
\end{enumerate}

\section{TECHNOLOGY STACK}

\subsection{Frontend: TypeScript SDK}

TypeScript was selected for the frontend SDK due to:
\begin{itemize}
    \item Static typing for reliability and maintainability
    \item Excellent IDE support and developer experience
    \item Seamless JavaScript interoperability
    \item Growing adoption in enterprise applications
\end{itemize}

The SDK provides:
\begin{itemize}
    \item Core \texttt{PassiveGuard} class for vanilla JavaScript
    \item \texttt{usePassiveGuard} React hook for React applications
    \item Event system for monitoring collection status
    \item Configurable privacy modes
\end{itemize}

\subsection{Backend: FastAPI}

FastAPI was chosen for the backend service due to:
\begin{itemize}
    \item Async support for high concurrency
    \item Automatic OpenAPI documentation
    \item Pydantic validation for request/response models
    \item Native Python type hints
    \item High performance (comparable to Node.js and Go)
\end{itemize}

\subsection{Machine Learning: XGBoost}

XGBoost (eXtreme Gradient Boosting) was selected for classification:
\begin{itemize}
    \item State-of-the-art performance on tabular data
    \item Fast training and inference (sub-50ms)
    \item Built-in feature importance analysis
    \item Robust handling of missing values
    \item Regularization to prevent overfitting
\end{itemize}

\subsection{Supporting Technologies}

\begin{table}[H]
\centering
\caption{Technology Stack Components}
\label{tab:tech_stack}
\begin{tabular}{|l|l|l|}
\hline
\textbf{Component} & \textbf{Technology} & \textbf{Purpose} \\
\hline
Frontend SDK & TypeScript 5.3+ & Type-safe data collection \\
\hline
React Integration & React 17+ & Hook-based integration \\
\hline
Backend API & FastAPI 0.109+ & REST API service \\
\hline
ML Framework & scikit-learn 1.4+ & Feature preprocessing \\
\hline
Classification & XGBoost 2.0+ & Bot detection model \\
\hline
Data Validation & Pydantic 2.5+ & Request/response schemas \\
\hline
Serialization & joblib 1.3+ & Model persistence \\
\hline
\end{tabular}
\end{table}

\section{DATA COLLECTION METHODOLOGY}

\subsection{Environmental Data Collection}

Environmental data captures static characteristics of the browser and device:

\begin{verbatim}
Environmental Features:
├── Browser Information
│   ├── User Agent String
│   ├── Platform
│   ├── Language
│   ├── Cookie Enabled
│   └── Do Not Track
├── Screen Properties
│   ├── Width, Height
│   ├── Color Depth
│   ├── Pixel Ratio
│   └── Orientation
├── Hardware Signals
│   ├── Device Memory
│   ├── Hardware Concurrency (CPU cores)
│   └── Max Touch Points
├── Feature Detection
│   ├── WebGL Support
│   ├── WebRTC Support
│   ├── Canvas Support
│   └── Audio API Support
└── Fingerprints (Hashed)
    ├── Canvas Fingerprint
    ├── WebGL Fingerprint
    └── Audio Fingerprint
\end{verbatim}

\subsection{Behavioral Data Collection}

Behavioral data captures dynamic user interaction patterns:

\subsubsection{Mouse Movement Analysis}

Mouse movements are captured at 60Hz and analyzed for:
\begin{itemize}
    \item \textbf{Velocity Statistics:} Mean, standard deviation, maximum velocity
    \item \textbf{Trajectory Analysis:} Straight-line ratio, curvature
    \item \textbf{Direction Changes:} Frequency of direction reversals
    \item \textbf{Pause Patterns:} Number and duration of pauses
    \item \textbf{Jitter Score:} High-frequency movement noise
\end{itemize}

The velocity between consecutive points is calculated as:

$$v_i = \frac{\sqrt{(x_{i+1} - x_i)^2 + (y_{i+1} - y_i)^2}}{t_{i+1} - t_i}$$

\subsubsection{Keyboard Pattern Analysis}

Keyboard interactions are analyzed for:
\begin{itemize}
    \item \textbf{Dwell Time:} Duration of key press (key down to key up)
    \item \textbf{Flight Time:} Interval between consecutive key presses
    \item \textbf{Typing Speed:} Characters per minute
    \item \textbf{Error Rate:} Backspace frequency
    \item \textbf{Rhythm Consistency:} Standard deviation of inter-key intervals
\end{itemize}

\subsubsection{Scroll Behavior Analysis}

Scroll events are analyzed for:
\begin{itemize}
    \item \textbf{Total Scroll Distance:} Cumulative pixels scrolled
    \item \textbf{Scroll Velocity:} Speed of scrolling
    \item \textbf{Smoothness:} Consistency of scroll increments
    \item \textbf{Direction Changes:} Frequency of scroll direction reversals
\end{itemize}

Scroll smoothness is calculated as:

$$\text{Smoothness} = 1 - \frac{\sigma(\Delta s)}{\mu(\Delta s)}$$

where $\Delta s$ represents scroll increments between events.

\subsubsection{Timing Analysis}

Temporal patterns provide additional signals:
\begin{itemize}
    \item \textbf{Time to First Interaction:} Delay before first user action
    \item \textbf{Total Interaction Time:} Duration of active engagement
    \item \textbf{Idle Periods:} Gaps between interactions
    \item \textbf{Interaction Gaps:} Distribution of inter-action intervals
\end{itemize}

\section{FEATURE ENGINEERING}

\subsection{Feature Vector Construction}

The raw collected data is transformed into a 56-dimensional feature vector:

\begin{table}[H]
\centering
\caption{Feature Categories}
\label{tab:features}
\begin{tabular}{|l|c|l|}
\hline
\textbf{Category} & \textbf{Count} & \textbf{Examples} \\
\hline
Environmental & 18 & Screen ratio, CPU cores, memory \\
\hline
Mouse Movement & 12 & Velocity mean/std, direction changes \\
\hline
Keyboard Patterns & 8 & Dwell time, flight time, speed \\
\hline
Scroll Behavior & 6 & Smoothness, velocity, distance \\
\hline
Timing & 8 & First interaction, total time, gaps \\
\hline
Fingerprints & 4 & Canvas, WebGL, audio (binary presence) \\
\hline
\textbf{Total} & \textbf{56} & \\
\hline
\end{tabular}
\end{table}

\subsection{Feature Normalization}

Numerical features are normalized using StandardScaler:

$$x_{normalized} = \frac{x - \mu}{\sigma}$$

where $\mu$ is the feature mean and $\sigma$ is the standard deviation from training data.

\subsection{Bot Indicator Features}

Specific features indicate potential bot characteristics:
\begin{itemize}
    \item \textbf{WebDriver Detection:} Presence of \texttt{navigator.webdriver} flag
    \item \textbf{Automation Flags:} Headless browser indicators
    \item \textbf{Plugin Anomalies:} Missing or unusual plugin configurations
    \item \textbf{Zero Interaction:} No mouse/keyboard events captured
\end{itemize}

\section{MACHINE LEARNING MODEL}

\subsection{XGBoost Classifier}

XGBoost implements gradient boosted decision trees with the following objective function:

$$\mathcal{L}^{(t)} = \sum_{i=1}^{n} l(y_i, \hat{y}_i^{(t-1)} + f_t(x_i)) + \Omega(f_t)$$

where $f_t$ is the tree added at iteration $t$, and $\Omega$ is the regularization term:

$$\Omega(f) = \gamma T + \frac{1}{2}\lambda \sum_{j=1}^{T} w_j^2$$

where $T$ is the number of leaves and $w_j$ are leaf weights.

\subsection{Model Hyperparameters}

\begin{table}[H]
\centering
\caption{XGBoost Hyperparameters}
\label{tab:hyperparams}
\begin{tabular}{|l|c|l|}
\hline
\textbf{Parameter} & \textbf{Value} & \textbf{Purpose} \\
\hline
n\_estimators & 100 & Number of boosting rounds \\
\hline
max\_depth & 6 & Maximum tree depth \\
\hline
learning\_rate & 0.1 & Step size shrinkage \\
\hline
subsample & 0.8 & Row sampling ratio \\
\hline
colsample\_bytree & 0.8 & Column sampling ratio \\
\hline
min\_child\_weight & 1 & Minimum leaf weight \\
\hline
reg\_alpha & 0 & L1 regularization \\
\hline
reg\_lambda & 1 & L2 regularization \\
\hline
\end{tabular}
\end{table}

\subsection{Training Procedure}

The model training procedure:
\begin{enumerate}
    \item Generate synthetic training data simulating human and bot behaviors
    \item Split data into training (80\%) and validation (20\%) sets
    \item Train XGBoost with early stopping on validation loss
    \item Evaluate using cross-validation (5-fold)
    \item Export trained model using joblib serialization
\end{enumerate}

\subsection{Prediction Scoring}

The model outputs a probability score $P(\text{human}|x) \in [0, 1]$. The verification decision follows:

$$\text{Decision} = \begin{cases} 
\text{HUMAN} & \text{if } P(\text{human}|x) \geq 0.7 \\
\text{SUSPICIOUS} & \text{if } 0.3 \leq P(\text{human}|x) < 0.7 \\
\text{BOT} & \text{if } P(\text{human}|x) < 0.3
\end{cases}$$

\section{API DESIGN}

\subsection{Verification Endpoint}

The primary verification endpoint accepts collected data and returns results:

\begin{verbatim}
POST /api/v1/verify
Content-Type: application/json

Request:
{
  "siteKey": "site_xxx",
  "environmental": { ... },
  "behavioral": { ... },
  "requestId": "pg_xxx",
  "timestamp": 1706620800000
}

Response:
{
  "isHuman": true,
  "confidence": 0.94,
  "riskScore": 0.12,
  "token": "eyJhbGc...",
  "challengeRequired": false
}
\end{verbatim}

\subsection{Token Validation}

Verification tokens can be validated server-side:

\begin{verbatim}
POST /api/v1/validate
Content-Type: application/json

Request:
{
  "token": "eyJhbGc...",
  "siteKey": "site_xxx"
}

Response:
{
  "valid": true,
  "isHuman": true,
  "confidence": 0.94,
  "timestamp": 1706620800000
}
\end{verbatim}

\section{PRIVACY CONSIDERATIONS}

\subsection{Data Minimization}

The system collects only data necessary for bot detection:
\begin{itemize}
    \item No personal identifiers (name, email, IP address in features)
    \item No persistent tracking across sessions
    \item No third-party data sharing
\end{itemize}

\subsection{Fingerprint Hashing}

All fingerprints are hashed using SHA-256 before transmission:

$$\text{hash} = \text{SHA256}(\text{fingerprint\_data})$$

This enables uniqueness comparison without storing raw fingerprint data.

\subsection{Privacy Mode}

An optional privacy mode disables fingerprinting while maintaining behavioral analysis, accepting slightly reduced accuracy for enhanced privacy.

